\section{Information and informing are different causal relations}
\label{sec:cbrd-information}

A difference may change a system's later action without being communicated. A message may be delivered without becoming consequential for its receiver. CBRD therefore separates:
\begin{align}
Information(\Delta;S)
&:= \Delta\text{ can originate or change a later transformation of }S,\\
Inform(i,j,\Delta)
&:= center\ i\text{ acts so that }\Delta\text{ is presented to center }j.
\end{align}
Hence
\begin{equation}
information\neq message,\qquad
delivery\neq uptake\neq retained\ consequence,
\end{equation}
and
\begin{equation}
informing(j)\neq authorship\text{-}over(j).
\end{equation}
Bateson's cybernetic formulation is a historical-priority neighbor for consequential information \cite{Bateson1972}. The project-specific residual, if any, lies in combining consequential difference with source/exposure provenance, participant jurisdiction, and future-use effects.

The terminology itself is part of the causal system. A sign may change which associations a receiver can access. This is why first reference must expose the referent and relation before a loaded project name is allowed to compress them.

\section{Typed claims before verdicts}

A claim status is not a property of a sentence detached from its research state. Define
\begin{equation}
\mathrm{TypedClaim}(Q)=\langle r,F,g,i,j,p,a,x,b,d\rangle,
\end{equation}
where the coordinates denote referent, frame, register, index/scale, jurisdiction, provenance, assumptions, realization/test object, dependency bridges, and discriminator. A material change in one of these coordinates requires status re-audit.

Three ternary structures must remain separate.

\paragraph{Developmental/root.}
\begin{equation}
|V\rangle\;|\Delta\rangle\;|T\rangle
\end{equation}
tracks a model from pre-stabilized relation potential through consequential differentiation to stabilized distinction.

\paragraph{Epistemic warrant.}
\begin{equation}
\mathrm{LICENSE}_F(P)\;|\;\mathrm{LICENSE}_F(\neg P)\;|\;\mathrm{UNDETERMINED}_F(P).
\end{equation}

\paragraph{Governance.}
\begin{equation}
\mathrm{ADMIT}\;|\;\mathrm{WITHHOLD}\;|\;\mathrm{REJECT/CUT}.
\end{equation}
An unresolved proposition does not mechanically dictate one governance action. Safety, jurisdiction, reversibility, and asymmetric cost may independently constrain action.

\section{Provenance as operative state}

When source or exposure can alter future consequence, provenance is not optional metadata. Represent state as
\begin{equation}
z=(\text{content},\text{source},\text{history},\text{exposure},\text{modality},\text{jurisdiction}).
\end{equation}
Dropping a coordinate requires a demonstrated invariance of the downstream transformations relevant to the claim.

The point is familiar from scientific provenance standards and data lineage, but CBRD pushes the requirement into recursive decision systems where the model may act on its own compression \cite{W3CProvDM2013}. Two otherwise identical current representations can have different future consequences if their sources trigger different trust, authorization, policy, or correction pathways.

\begin{definition}[Source-swap test]
Replace the source while preserving the apparent content. If the admissible action, expected consequence, or verification requirement changes, source is part of operative state for that task.
\end{definition}

\section{Recursive sufficiency}

Ordinary compression asks whether a representation preserves enough information for a present task. Recursive use asks a harder question: does the compression remain sufficient after the representation influences future action, policy, observation, or model state?

For a quotient $\phi:X\rightarrow X/{\sim}$, call $\phi$ recursively sufficient for a declared transformation family $\mathcal{T}$ and horizon $H$ only when distinctions erased by $\phi$ remain irrelevant under the future transformations induced through the quotient itself.

This connects to state abstraction, bisimulation, causal abstraction, computational mechanics, and task-relative macrostate construction \cite{GivanDeanGreig2003,LiWalshLittman2006,HoelEtAl2013,Hoel2017,GeigerEtAl2021,GeigerEtAl2023}. The residual claim is not that these traditions missed compression. It is that self-conditioning, provenance activation, policy change, and returned consequence can make a currently irrelevant distinction matter later.

A high-resolution model can become psychologically harder to falsify while still omitting the same bridge. Resolution therefore cannot be used as a substitute for explicit residuals.

\section{Returned consequence and claim-relative non-preauthorship}

A system that acts on its own predictions contaminates the later observation distribution. This does not make external correction impossible; it makes independence claim-relative.

Let $R$ be a returned observation and $H$ the focal hypothesis or policy. The useful question is not whether $R$ is causally pure. It is whether enough test-relevant variation in $R$ was outside $H$'s prior control to discriminate the claim under review.

\begin{definition}[Claim-relative causal non-preauthorship]
A returned consequence is sufficiently non-preauthored for claim $Q$ when the variation used to update $Q$ was not fully fixed by the model, policy, or evidence route whose adequacy $Q$ asserts.
\end{definition}

Thus
\begin{align}
\text{feedback}&\neq\text{independent correction},\nonumber\\
\text{surprise}&\neq\text{truth},\nonumber\\
\text{source multiplicity}&\neq\text{causal independence}.
\end{align}
Performative prediction and feedback-loop literatures are relevant controls here \cite{PerdomoEtAl2020}.

\section{Consequence-bearing recursion and exposure}
\label{sec:deep-recursion}

The preceding sections can be used without a special theory of recursion. The deeper research program, however, needs to distinguish three superficially similar loops before causal claims are made.

\begin{definition}[Mere recurrence]
A process exhibits mere recurrence when a transformation repeatedly maps prior state into later state, for example $x_{t+1}=F(x_t)$, without requiring returned consequence to revise the transformation itself.
\end{definition}

\begin{definition}[Self-conditioned completion]
A process exhibits self-conditioned completion when its own outputs materially shape later inputs, labels, observations, or evaluation conditions. Internal coherence can then increase while exposure to independently originating correction decreases.
\end{definition}

\begin{definition}[Consequence-bearing recursion]
A process exhibits consequence-bearing recursion when an emitted action or state reaches an environment or another center, a returned consequence is not fully pre-authored by the focal process for the tested claim, relevant provenance is retained, and that consequence can change later model, policy, interface, generator, or authorization state.
\end{definition}

The point is not that independent evidence is metaphysically pure. The useful variable is claim-relative causal dependence. Let
\begin{equation}
\mathcal X^{\xi}_{i\rightarrow j}(t)
\label{eq:deep-exposure}
\end{equation}
index the exposure of returning source $j$ to source $i$'s framing, desired conclusion, intervention, or prediction about task variable $\xi$. Mature applications should replace this placeholder by an explicit causal exposure graph. A returned statement after heavy exposure can remain excellent evidence about the participant's current response while being weak evidence for independent generation of the hypothesis being tested.

A provenance class can therefore distinguish, at declared scope,
\begin{equation}
\operatorname{pclass}(e)\in
\{independent,exposed,model\!\text{-}guided,self\!\text{-}generated,mixed,unknown\}.
\label{eq:deep-pclass}
\end{equation}
These labels are not a universal ranking. They are variables whose usefulness must itself be tested.

\paragraph{A local informational now.}
For boundary $B_i$, pre-contact state $s_i^-$, and incoming difference $\Delta$, define
\begin{equation}
Now_{B_i}(\Delta)
\Longleftrightarrow
\Delta\text{ crosses }B_i\text{ and }s_i^+\neq s_i^-.
\label{eq:deep-now}
\end{equation}
This is an incidence definition, not a consciousness claim and not a derivation of fundamental physical time. Different centers can have different or partially ordered ``nows'' because consequential contact is indexed to their boundaries.

\begin{readercheck}
Remove the labels. Distinguish repeated internal processing, a loop whose later evidence is strongly shaped by its own earlier outputs, and a loop in which returned consequence can rewrite persistent future organization. If the distinction produces no new diagnostic or intervention in the target domain, ordinary feedback terminology is enough.
\end{readercheck}
